\documentclass[10pt,a4paper]{amsart}
\usepackage[margin=19mm]{geometry}
\usepackage{amsmath,amssymb,mathtools}
\usepackage[scr=rsfs,cal=euler]{mathalfa}
\usepackage{xfrac}
\usepackage[unicode,hidelinks,bookmarks=false]{hyperref}
\newcommand{\ie}{i.e.\ }
\newcommand{\cf}{cf.\ }
\newcommand{\resp}{respectively\ }
\newcommand{\Subsection}{\subsection}
\providecommand{\nobreakdash}{\nobreak\hbox{-}}
\setlength{\emergencystretch}{2em}
\multlinegap0pt
\newtheorem{proposition}{Proposition}
\title{Constructions of Frenet Curves with respect to Semi-Symmetric Metric Connection}
\author{\c{S}aban G\"uven\c{c}}
\date{Source: arXiv:2408.05880v1 (CC BY 4.0)}
\begin{document}
\maketitle

\noindent\emph{Source and licence.} Constructions of Frenet Curves with respect to Semi-Symmetric Metric Connection. Version arXiv:2408.05880v1 (CC BY 4.0), distributed by arXiv under the Creative Commons Attribution 4.0 International licence, http://creativecommons.org/licenses/by/4.0/, as recorded in the arXiv licence metadata for this exact version and shown on the abstract page https://arxiv.org/abs/2408.05880v1. Full licence notice: \texttt{licenses/arxiv-2408.05880-CC-BY-4.0.txt}.\par\smallskip

\noindent\textit{Editorial scope.} Counted unit: the article's proposition proposition (source0.tex lines 205--210). The article's complete original proof of it is delivered with it (source0.tex lines 212--251, one \texttt{proof} environment of 40 lines). No mathematics is written, repaired or re-derived: the statement and the proof are the article's own lines, in the article's own notation, and no retained line is substituted or re-flowed.  No further source-local context is needed: every cross-reference of every retained line resolves inside the two delivered spans.  Nothing was omitted from any retained span, so the package carries no omission record.  No retained line contains a citation, so the wrapper carries no bibliography and no bibliography entry.  No retained line contains a figure environment, an image-inclusion command or drawing code, so no image is delivered and the package declares no asset dependency.  Editorial formatting only, in addition: an AMS \texttt{amsart} preamble that restates the article's own declarations and macros used by the retained lines, namely the environments \texttt{proposition} on separate counters, together with the standard packages those lines need.  The retained environments print in this document as Proposition 1 in order of appearance, whereas the article prints the same items with the numbering of its own full document.  The complete native source of the article as distributed in the arXiv e-print for this exact version is bundled unchanged as \texttt{sources/163655/source0.tex}.
\begin{proposition}
For a semi-symmetric metric Frenet curve of order $3$, 
\begin{equation*}
\widetilde{\nabla }_{T}B=-\widetilde{\tau }N.
\end{equation*}
\end{proposition}

\begin{proof}
Since $M$ is $3$-dimensional, we $\widetilde{\nabla }_{T}B\in span\left\{
T,N,B\right\} $, which is a $g$-orthonormal vector basis along the curve $%
\gamma $ by construction. So we can write%
\begin{equation}
\widetilde{\nabla }_{T}B=g\left( \widetilde{\nabla }_{T}B,T\right) T+g\left( 
\widetilde{\nabla }_{T}B,N\right) N+g\left( \widetilde{\nabla }%
_{T}B,B\right) B.  \label{0}
\end{equation}%
Differentiating $g\left( T,B\right) =0,$%
\begin{equation*}
g\left( \widetilde{\nabla }_{T}T,B\right) +g\left( T,\widetilde{\nabla }%
_{T}B\right) =0,
\end{equation*}%
\begin{equation*}
g\left( \widetilde{\kappa }N,B\right) +g\left( T,\widetilde{\nabla }%
_{T}B\right) =0,
\end{equation*}%
\begin{equation}
g\left( T,\widetilde{\nabla }_{T}B\right) =0.  \label{1}
\end{equation}%
Differentiating $g\left( N,B\right) =0,$%
\begin{equation*}
g\left( \widetilde{\nabla }_{T}N,B\right) +g\left( N,\widetilde{\nabla }%
_{T}B\right) =0,
\end{equation*}%
\begin{equation*}
g\left( -\widetilde{\kappa }T+\widetilde{\tau }B,B\right) +g\left( N,%
\widetilde{\nabla }_{T}B\right) =0,
\end{equation*}%
\begin{equation}
g\left( N,\widetilde{\nabla }_{T}B\right) =-\widetilde{\tau }.  \label{2}
\end{equation}%
Differentiating $g\left( B,B\right) =1,$%
\begin{equation}
g\left( \widetilde{\nabla }_{T}B,B\right) =0.  \label{3}
\end{equation}%
If we write (\ref{1}),(\ref{2}) and (\ref{3}) in (\ref{0}), we get $%
\widetilde{\nabla }_{T}B=-\widetilde{\tau }N.$
\end{proof}

\end{document}
